\subsection{媒介変数表示}
曲線が $x=x(t)$,\,$y=y(t)$ と表されるとき,\,$t$ を\textbf{媒介変数}という.ここでは $x(t),y(t)$ が連続な導関数をもつとする.

考えている点で $x'(t)\ne0$ なら,\,$x'$ の連続性からその近くでも $x'(t)$ は同じ符号をもち,\,$x(t)$ は狭義単調になる.したがって,前の小節により $t$ を $x$ の関数 $t=t(x)$ として表せて,
\[
 \frac{\dd t}{\dd x}=\frac1{x'(t)}
\]
となる.\,$y=y(t(x))$ は $x$ の関数だから,合成関数の微分により
\begin{align*}
 \frac{\dd y}{\dd x}
 &=y'(t)\frac{\dd t}{\dd x}\\
 &=\frac{y'(t)}{x'(t)}
\end{align*}
を得る.\,$x'(t)=0$ の点ではこの商は使えず,垂直な接線などを別に調べる必要がある.

さらに2回微分できる場合は,同じ合成関数の微分から,\,$x$ で微分することが $\frac1{x'(t)}\frac{\dd}{\dd t}$ に対応する.よって $\frac{\dd y}{\dd x}$ をもう一度 $x$ で微分すると
\[
 \frac{\dd^2y}{\dd x^2}
 =\frac1{x'(t)}\frac{\dd}{\dd t}
   \left(\frac{y'(t)}{x'(t)}\right)
\]
となる.\,$\frac{y''(t)}{x''(t)}$ ではないことに注意しよう.
